
\documentclass[10pt, a4paper]{article}

% --- PACKAGES ---
\usepackage[utf8]{inputenc}
\usepackage[T1]{fontenc}
\usepackage[margin=1.2cm]{geometry}
\usepackage{hyperref}
\usepackage{enumitem}
\usepackage[explicit]{titlesec}
\usepackage{xcolor}
\usepackage{fontawesome5}
\usepackage{lmodern}

% --- COLOR DEFINITIONS ---
\definecolor{primary}{RGB}{30, 50, 90}      % Navy Blue accent
\definecolor{textdark}{RGB}{40, 40, 40}    % Dark Charcoal for text
\definecolor{subtle}{RGB}{100, 100, 100}  % Grey for muted text

% --- SETUP & STYLING ---
\pagestyle{empty}
\color{textdark}
\hypersetup{
    colorlinks=true,
    linkcolor=primary,
    urlcolor=primary
}

% Custom section formatting
\titleformat{\section}
  {\color{primary}\large\bfseries}
  {}{0em}
  {\MakeUppercase{#1}}
  [\color{primary}\titlerule]

\titlespacing*{\section}{0pt}{8pt}{4pt}

% Custom list parameters
\setlist[itemize]{leftmargin=*, parsep=0pt, itemsep=2pt, topsep=1pt}

% --- DOCUMENT BEGIN ---
\begin{document}

% --- HEADER ---
\begin{center}
    {\Huge \bfseries \color{primary} Jean MARTINI}\\[4pt]
    {\large \bfseries Ingénieur IA -- Diplômé de Centrale Lyon}\\[6pt]
    \small
    Paris \quad$\cdot$\quad 
    +33 7 52 02 31 20 \quad$\cdot$\quad 
    \href{mailto:jean.martini8900@gmail.com}{jean.martini8900@gmail.com} \quad$\cdot$\quad 
    Permis B\\[2pt]
    \href{https://linkedin.com/in/jean-martini-64a017228/}{\faLinkedin\ LinkedIn} \quad$\cdot$\quad 
    \href{https://github.com/Jean8900}{\faGithub\ GitHub}
\end{center}

% --- PROFILE ---
\section{Profil}
Diplômé de Centrale Lyon (2025), spécialisé en modélisation prédictive, vision par ordinateur et grands modèles de langage (LLM). Participant actif à des hackathons d'IA (1\,750 € remportés). Actuellement Data Scientist / Ingénieur IA chez Illuin Technology.\\[3pt]
\textbf{Tech Stack:} Python, PyTorch, scikit-learn, XGBoost, OpenCV, Flask, Docker, SQL, Plotly, Git, Mistral API\\
\textbf{Langues :} Français (langue maternelle) \quad$\cdot$\quad Anglais (C1) \quad$\cdot$\quad Espagnol (B1) \quad$\cdot$\quad Japonais (A2)

% --- WORK EXPERIENCE ---
\section{Expérience professionnelle}

\textbf{Illuin Technology} \hfill Paris, France\\
\textit{Ingénieur IA} \hfill \small Avril 2026 -- Aujourd'hui
\begin{itemize}
    \item Développement et amélioration d'agents conversationnels avec ILLUIN Dialogue pour les parcours de souscription de forfaits mobiles et de gestion de compte d'Orange.
    \item Conception de 5 agents spécialisés mobilisant prompt engineering, tool calling, RAG, tests automatisés et optimisation des parcours conversationnels.
    \item Travail au sein d'une équipe Illuin de 20 personnes, en lien étroit avec les équipes d'Orange, pour garantir la fiabilité des réponses et des actions des agents.
\end{itemize}

\textbf{IFP Energies Nouvelles} \hfill Lyon, France\\
\textit{Stagiaire Data Scientist} \hfill \small Avril 2025 -- Septembre 2025
\begin{itemize}
    \item Classification industrielle de pneus par scan spectral pour le recyclage : maîtrise de la technologie de scan NIR, protocole expérimental, acquisition des données et conception du pipeline ML (PLS-DA + SVM).
    \item Mise en œuvre de l'algorithme des projections successives (SPA) pour la sélection de variables : temps d'acquisition divisé par 10 avec 100\% de précision conservée.
    \item Poster présenté au Congrès national de chimiométrie (GFC-NCA 2026, février 2026).
\end{itemize}

\textbf{Société du Canal de Provence} \hfill Aix-en-Provence, France\\
\textit{Stagiaire / Prestataire analyste données climatiques} \hfill \small Septembre 2023 -- Mars 2024
\begin{itemize}
    \item Développement de pipelines d'analyse climatique : intégration de jeux de données hydro-climatiques mondiaux (WorldClim, DRIAS, Explore2) et de projections du changement climatique à l'horizon 2050 (scénarios SSP/RCP).
    \item Réalisation d'études prédictives internationales sur la dynamique des nappes phréatiques à Madagascar, les risques de sécheresse au Panama et l'hydroélectricité en Bosnie (capacité projetée -20\% d'ici 2050). Collaboration avec Ernst \& Young. Stage prolongé de 3 mois.
\end{itemize}

% --- EDUCATION ---
\section{Formation}

\textbf{Centrale Lyon} \hfill Septembre 2021 -- Septembre 2025\\
\textit{Diplôme d'ingénieur généraliste (Master)} \hfill Lyon, France\\[2pt]
\textbf{Keio University} \hfill Avril 2023 -- Août 2023\\
\textit{Échange académique en intelligence artificielle} -- spécialisation Data Science \hfill Tokyo, Japon\\
\small Échange académique centré sur des cours d'IA et de vision par ordinateur.

\vspace{2pt}
\noindent\textbf{Certifications :} Dataiku Core Designer (2025) \quad$\cdot$\quad Stanford Deep Learning (Andrew Ng) \quad$\cdot$\quad Inria Machine Learning with scikit-learn

% --- HACKATHONS & COMPETITIONS ---
\section{Hackathons \& Compétitions}

\textbf{GE Healthcare - Mistral AI - Hackathon AI for Health} \hfill \textbf{3\textsuperscript{e} / 22 équipes (prix de 1\,750 €)}
\begin{itemize}
    \item Segmentation par vision par ordinateur de scanners CT 3D : extraction automatique de la coupe montrant le diamètre maximal de la tumeur (2\,000 vues).
    \item Classification Lung-RADS / RECIST et génération automatique de comptes rendus via Mistral AI, validées par un radiologue.
    \item Temps d'analyse divisé par 10 (20 min $\rightarrow$ 2 min, réf. IGAS) pour un coût de 0,01 €/compte rendu sur données réelles.
\end{itemize}

\textbf{CSIRO Image2biomass} (cagnotte de 75\,k\$) \hfill \textbf{Top 10\%}
\begin{itemize}
    \item Compétition de vision par ordinateur : prédiction de la biomasse agricole à partir d'images de terrain.
    \item Combinaison Vision Transformer (DINOv2) + MLP pour atteindre $R^2 = 0.70$.
\end{itemize}

% --- VOLUNTEERING & LEADERSHIP ---
\section{Engagement associatif \& Leadership}

\textbf{Adopt AI Forum} \hfill Novembre 2025\\
\textit{Forum mondial sur l'adoption de l'IA dans les industries (600 PDG, keynote du Président de la République)}
\begin{itemize}
    \item Responsable des espaces privés : coordination de la logistique et interlocuteur principal des partenaires entreprises (BNP Paribas, Nvidia, Meta, Servier).
\end{itemize}

\textbf{Solidari'terre} \hfill 2022 -- 2023\\
\textit{Chef de projet -- Maraudes de nuit (association humanitaire de Centrale Lyon, 17 membres)}
\begin{itemize}
    \item Coordination de 5 bénévoles par session, gestion des stocks et de la logistique pour 30 distributions de repas par semaine.
\end{itemize}

\end{document}

